\section{Metode zasnovane na unutrašnjoj strukturi modela}

Metode interpretabilnosti veštačke inteligencije (\textit{Explainable AI} - XAI) mogu se fundamentalno podeliti na osnovu stepena pristupa unutrašnjoj arhitekturi i parametrima modela. Ovo poglavlje posvećeno je metodama specifičnih za model ( \textit{white-box methods}) koje zahtevaju potpun uvid u unutrašnje stanje mreže, uključujući njene težine, aktivacije i gradijente. Nasuprot njima, metode nezavisne od modela, koje sistem posmatraju isključivo spolja, biće detaljno analizirane u narednom poglavlju.

Unutar ovog poglavlja, fokus je usmeren na tri ključne grupe pristupa. Najpre se analiziraju istorijski i konceptualno fundamentalne metode zasnovane na gradijentima. Iako ove metode ne predstavljaju primarni fokus evaluacije u ovom radu, one služe kao neophodna teorijska osnova i intuitivna motivacija za razumevanje naprednijih pristupa. Analiza njihovih prednosti, a naročito ograničenja, direktno otvara prostor za sofisticiranije metode specifične za arhitekture zasnovane na mehanizmu pažnje. U tom kontekstu, centralni deo poglavlja posvećen je metodama propagacije relevantnosti po slojevima (LRP) i tehnikama koje eksploatišu samu strukturu Transformer modela, kao što su Attention Rollout \cite{abnar2020quantifying} i metoda propagacije relevantnosti kroz mehanizam pažnje \cite{chefer2021transformer}.

\subsection{Metode zasnovane na gradijentima}

Gradijentne metode interpretabilnosti zasnivaju se na ideji da parcijalni izvodi 
izlaza modela po ulaznim pikselima nose informaciju o tome koliko svaki piksel 
doprinosi odluci modela. Neka je $f_c(\mathbf{x})$ izlaz modela za klasu $c$ pri 
ulaznoj slici $\mathbf{x} \in \mathbb{R}^{H \times W \times C}$. Gradijent izlaza po ulaznoj slici definiše se kao:

\begin{equation}
    \mathbf{G} = \frac{\partial f_c(\mathbf{x})}{\partial \mathbf{x}}
\end{equation}

Dobijeni tenzor $\mathbf{G}$ zapravo pokazuje koliko je model osetljiv na promenu svakog pojedinačnog piksela. Ako mala promena vrednosti nekog piksela izaziva veliku promenu u izlazu modela, gradijent za taj piksel će biti visok. Na taj način, gradijent direktno otkriva koje delove slike model smatra najvažnijim za donošenje konačne odluke.

\subsubsection{Vanilla Gradient}

Vanilla Gradient \cite{simonyan2014deep} predstavlja najjednostavniji pristup 
interpretabilnosti zasnovan na gradijentima. Mapa značajnosti dobija se direktno kao 
apsolutna vrednost gradijenta izlaza modela po ulaznoj slici:

\begin{equation}
    S_{vanilla}(\mathbf{x}) = \left|\frac{\partial f_c(\mathbf{x})}{\partial \mathbf{x}}\right|
\end{equation}

Intuicija iza ove metode zasniva se na pretpostavci da piksel čija minimalna promena 
vrednosti dovodi do velike varijacije na izlazu modela može se smatrati ključnim za 
donošenje odluke. Mapa značajnosti generiše se kroz jedan prolaz unazad 
(\textit{backpropagation}) kroz mrežu, što ovu metodu čini računarski izuzetno efikasnom.

Uprkos svojoj jednostavnosti, Vanilla Gradient pati od dva fundamentalna 
ograničenja koja ozbiljno narušavaju pouzdanost dobijenih objašnjenja.

Prvo ograničenje odnosi se na problem zasićenja (\textit{saturation problem}). Budući da gradijent opisuje isključivo lokalnu osetljivost modela u specifičnoj tački $\mathbf{x}$, karakteristike koje su suštinski ključne za donošenje odluke mogu imati gradijent blizak nuli ukoliko se model nalazi u zoni zasićenja aktivacionih funkcija. Ovaj fenomen nastaje kada su određena vizuelna obeležja već dovoljno izražena na ulazu da u potpunosti aktiviraju neurone odgovorne za prepoznavanje posmatrane klase. U takvoj situaciji, dodatno povećanje njihovog intenziteta ne utiče na izlaz modela, pa gradijent postaje veoma mali ili jednak nuli. Kao posledica toga, karakteristike koje nose ključnu diskriminativnu informaciju mogu biti pogrešno ocenjene kao nerelevantne.

Drugo ograničenje tiče se vizuelne nestabilnosti dobijenih mapa. Usled 
višestrukih nelinearnih transformacija kroz duboke slojeve mreže, gradijent postaje 
izrazito nestabilan i podložan naglim lokalnim fluktuacijama čak i za međusobno 
veoma slične ulazne piksele — fenomen poznat kao problem drobljenja gradijenata 
(\textit{shattered gradients problem}) \cite{balduzzi2017shattered}. Posledica je 
vizuelno nejasna, tačkasta mapa značajnosti koja otežava jasnu interpretaciju, 
što direktno motiviše unapređenja opisana u nastavku.

\subsubsection{Input $\times$ Gradient}

Kako bi se ublažio problem vizuelne nestabilnosti, predložena je metoda \textit{Input $\times$ Gradient} \cite{kindermans2016investigating}. Ovaj pristup linearno kombinuje prethodno definisani gradijent sa samom vrednošću ulaznog piksela:

\begin{equation}
    S_{I \times G}(\mathbf{x}) = \mathbf{x} \odot \frac{\partial f_c(\mathbf{x})}{\partial \mathbf{x}}
\end{equation}

gde $\odot$ označava Adamarov (element-po-element) proizvod. Teorijska 
opravdanost ovog pristupa proizilazi iz njegove interpretacije kao prvog člana 
Tejlorovog reda funkcije $f_c$ razvijenog oko nule \cite{bach2015pixel}:

\begin{equation}
    f_c(\mathbf{x}) \approx f_c(\mathbf{0}) + \nabla f_c(\mathbf{x}) \cdot \mathbf{x} + \mathcal{O}(\|\mathbf{x}\|^2)
\end{equation}

Član $\nabla f_c(\mathbf{x}) \cdot \mathbf{x}$ aproksimira doprinos svakog piksela 
izlazu modela pod pretpostavkom da je gradijent konstantan duž linearne putanje od 
nule do trenutne vrednosti piksela. Piksel koji ima istovremeno veliku vrednost i 
veliki gradijent dobija visoku ocenu značajnosti — on i snažno aktivira model i 
model je osetljiv na njegovu promenu.

U poređenju sa Vanilla Gradient, ovaj pristup uspešno eliminiše lažne 
pozitivne rezultate u regionima gde informacija na slici faktički ne postoji. 
Na primer, ukoliko model pokazuje nasumičnu osetljivost na crvenu boju u delu pozadine gde te boje zapravo nema, množenje sa nultom vrednošću ulaza ispravno anulira taj artefakt.

Međutim, uvođenje ulaza u formulu ne rešava problem zasićenja. Pretpostavka o 
konstantnom gradijentu duž putanje od $\mathbf{0}$ do $\mathbf{x}$ retko je 
opravdana kod dubokih mreža sa složenim nelinearnostima — ukoliko je gradijent 
na ključnim karakteristikama već pao na nulu usled zasićenja, množenje sa vrednošću 
ulaza ne može kompenzovati izgubljenu informaciju.

\subsubsection{GradCAM}

\textit{Gradient-weighted Class Activation Mapping} (GradCAM) \cite{selvaraju2017gradcam} 
pomera fokus sa ulaznog prostora na unutrašnje slojeve mreže. Motivacija leži u tome da 
dublji slojevi nose semantički bogatije reprezentacije, a njihovi gradijenti su stabilniji 
jer nisu izloženi kumulativnom efektu nelinearnih transformacija.

\paragraph{Konvolucione neuralne mreže.}
Kod CNN arhitektura primena je direktna: konvolucioni slojevi očuvavaju lokalnu prostornu 
korespondenciju, što omogućava precizno mapiranje aktivacija na regione ulazne slike. 
GradCAM koristi mape aktivacija poslednjeg konvolucionog sloja $\mathbf{A}^k$ i računa 
težinu svakog kanala globalnim usrednjavanjem gradijenta ciljnog izlaza po prostornim 
dimenzijama:
\begin{equation}
    \alpha_k^c = \frac{1}{H' W'} \sum_{i,j} \frac{\partial f_c}{\partial A^k_{ij}}
\end{equation}
Mapa značajnosti se formira kao težinska suma po svim kanalima, propuštena kroz ReLU aktivacionu funkciju ---
negativni doprinosi, koji umanjuju verovatnoću klase $c$, se odbacuju:
\begin{equation}
    S_{\text{GradCAM}}(\mathbf{x}) = \max\!\left(\sum_k \alpha_k^c \mathbf{A}^k,\; 0\right)
\end{equation}

\paragraph{Primena na Transformer arhitekture.}
Prilagođavanje ViT modelima nije trivijalno jer ViT podatke procesira kao sekvencu tokena, 
bez eksplicitne prostorne strukture u unutrašnjim slojevima. Konceptualna paralela se 
uspostavlja između CNN kanala i komponenenata pažnje (\textit{multi-head attention}): za blok $l$ 
i glavu $h$, aktivacioni vektor se konstruiše iz reda matrice pažnje koji odgovara CLS 
tokenu, pri čemu se CLS element izostavlja jer ne nosi prostornu informaciju:
\begin{equation}
    \mathbf{a}^{(l,h)} = \mathbf{A}^{(l,h)}_{[\text{CLS},\, 1:N]} \in \mathbb{R}^{N}
\end{equation}
Težina svake glave dobija se usrednjavanjem gradijenata po tokenima, a konačna mapa 
analogno CNN slučaju:
\begin{equation}
    \alpha_h^c = \frac{1}{N} \sum_{i=1}^{N} \frac{\partial f_c}{\partial a^{(l,h)}_i}, 
    \qquad
    S_{\text{GradCAM}}^{\text{ViT}}(\mathbf{x}) = 
    \max\!\left(\sum_h \alpha_h^c \cdot \mathbf{a}^{(l,h)},\; 0\right)
\end{equation}
Rezultujući vektor od $N$ elemenata direktno odgovara $N$ pečeva na koje je ulazna slika 
podeljena, pri čemu svaki peč dobija jedinstvenu skalarnu težinu. Mapa se zatim 
reorganizuje u mrežu dimenzija $\frac{H}{P} \times \frac{W}{P}$ i skalira bilinearnom 
interpolacijom na originalnu rezoluciju. Rezolucija je dakle inherentno ograničena 
veličinom pečeva --- tipično $14 \times 14$ za ViT-B/16 --- što predstavlja fundamentalno 
ograničenje zajedničko svim metodama koje će se ubuduće razmatrati u ovom radu i neće se posebno isticati 
u svakom slučaju. U praksi se koristi poslednji Transformer blok 
\cite{chefer2021transformer} kao semantički najkonkretniji.


%\subsection{Dodatna unapređenja gradijentnih metoda}
%
%Pored opisanih osnovnih metoda, u literaturi su predložena i brojna unapređenja koja 
%adresiraju njihove specifične nedostatke. U ovom radu ova unapređenja nisu bila 
%predmet evaluacije, ali predstavljaju prirodan pravac daljeg istraživanja.

%\subsection{Napredna unapređenja gradijentnih metoda na nivou piksela}
%
%Kako bi se prevazišli fundamentalni nedostaci \textit{Vanilla Gradient} i %\textit{Input $\times$ Gradient} metoda, u literaturi su predloženi robusniji %matematički okviri. Iako ovi pristupi nisu primarni predmet empirijske evaluacije u %ovom radu, njihovo razumevanje je ključno za sagledavanje šireg konteksta razvoja %metoda interpretabilnosti.

%\textbf{Integrated Gradients} \cite{sundararajan2017axiomatic} na elegantan način %rešava problem zasićenja gradijenata. Umesto da se osetljivost modela procenjuje %isključivo u pojedinačnoj tački $\mathbf{x}$, ova metoda integriše gradijente duž %linearne putanje od definisane referentne (bazne) slike $\mathbf{x}'$ (što je tipično %potpuno crna slika) do originalnog ulaza $\mathbf{x}$:

%\begin{equation}
%    S_{IG}(\mathbf{x}) = (\mathbf{x} - \mathbf{x}') \odot \int_0^1 \frac{\partial %f_c(\mathbf{x}' + \alpha(\mathbf{x} - \mathbf{x}'))}{\partial \mathbf{x}} \, d\alpha
%    \label{eq:integrated_gradients}
%\end{equation}

%U praksi se ovaj integral aproksimira Riemann-ovom sumom kroz $m$ diskretnih koraka. %Iako uspešno rešava problem zasićenja i zadovoljava aksiom kompletnosti, ova %aproksimacija zahteva $m$ prolaza unazad (\textit{backpropagation}) kroz mrežu umesto %jednog, što metodu čini računarski znatno zahtevnijom.

%\textbf{SmoothGrad} \cite{smilkov2017smoothgrad} direktno ublažava problem vizuelne %nestabilnosti i šuma u mapama značajnosti. Tehnika funkcioniše tako što generiše $n$ %perturbovanih verzija ulazne slike dodavanjem kontrolisanog Gaussovog šuma, a zatim %usrednjava gradijente izračunate nad svim pojedinačnim uzorcima:
%
%\begin{equation}
%    S_{SG}(\mathbf{x}) = \frac{1}{n}\sum_{i=1}^{n} \frac{\partial f_c(\mathbf{x} + %\epsilon_i)}{\partial \mathbf{x}}, \quad \epsilon_i \sim \mathcal{N}(0, \sigma^2)
%    \label{eq:smoothgrad}
%\end{equation}
%
%Ključna prednost \textit{SmoothGrad} pristupa ogleda se u njegovoj univerzalnosti. On %se ne posmatra kao izolovana tehnika, već kao opšti dodatak koji se može kombinovati %sa drugim pristupima (npr. \textit{SmoothGrad $\times$ Input} ili \textit{Smooth %Integrated Gradients}), efikasno uklanjajući visokofrekventni šum iz bilo koje %gradijentne mape.


\subsection{Propagacija relevantnosti po slojevima (LRP)}

Metode zasnovane na gradijentima, iako računarski efikasne, pate od fundamentalnog 
problema — gradijent opisuje lokalnu osetljivost modela, a ne stvarni doprinos ulaznih 
karakteristika izlaznoj odluci. Propagacija relevantnosti po slojevima 
(\textit{Layer-wise Relevance Propagation}, LRP) \cite{bach2015pixel} pristupa problemu 
interpretabilnosti iz potpuno drugačijeg ugla: umesto da pita kako bi se izlaz 
promenio kada bi se perturbovao ulaz, LRP pita koliko je svaki neuron doprineo trenutnoj odluci modela.

\subsubsection{Konzervacija relevantnosti}

Centralni koncept LRP metode je \textbf{relevantnost} $R_j^{(l)}$ — skalar koji kvantifikuje doprinos neurona $j$ u sloju $l$ konačnoj odluci modela. Relevantnost zadovoljava dva fundamentalna svojstva. Prvo, ukupna količina relevantnosti se konzervira pri prelasku između slojeva tokom propagacije unazad:

\begin{equation}
    \sum_j R_j^{(l)} = \sum_i R_i^{(l+1)} = R_{\text{total}}
    \label{eq:conservation}
\end{equation}

Drugo, relevantnost svakog pojedinačnog čvora dobija se potpunom raspodelom vrednosti između susednih slojeva:

\begin{equation}
    R_j^{(l)} = \sum_i R_{j \leftarrow i}^{(l, l+1)} \quad \text{uz uslov} \quad R_i^{(l+1)} = \sum_j R_{j \leftarrow i}^{(l, l+1)}
\end{equation}

gde $R_{j \leftarrow i}^{(l, l+1)}$ označava količinu relevantnosti koju čvor $i$ iz višeg sloja $l+1$ prosleđuje čvoru $j$ u nižem sloju $l$. Intuitivno, relevantnost se ponaša u skladu sa zakonom održanja — ona se ne stvara i ne gubi, već se isključivo redistribuira sa izlaznih prema ulaznim slojevima mreže.

Propagacija se inicijalizuje na poslednjem izlaznom sloju $L$ pre aktivacione funkcije Softmax. Za izabranu ciljnu klasu $c$, početna relevantnost se postavlja na vrednost mrežne aktivacije $f_c(\mathbf{x})$, dok se za sve ostale klase postavlja nula:

\begin{equation}
    R_k^{(L)} = 
    \begin{cases}
        f_c(\mathbf{x}) & \text{za } k = c \\
        0 & \text{za } k \neq c
    \end{cases}
\end{equation}
\subsubsection{Duboka Tejlorova dekompozicija (DTD) kao teorijska osnova}

Teorijsko utemeljenje LRP pravila pruža duboka Tejlorova dekompozicija (\textit{Deep Taylor Decomposition, DTD}) 
\cite{montavon2017explaining}. Osnovna ideja je da se relevantnost propagira 
unazad kroz mrežu sloj po sloj — analogno propagaciji gradijenata, ali sa 
drugačijom semantikom: umesto osetljivosti, propagira se doprinos.

Za jedan sloj sa operacijom $L(\mathbf{u}, \mathbf{v})$, gde su $\mathbf{u}$ 
ulazni tenzor a $\mathbf{v}$ tenzor parametara (ili drugi ulazni tenzor u 
slučaju neparametarskih operacija), Tejlorov razvoj oko korene tačke 
$\tilde{\mathbf{u}}$ takve da $L(\tilde{\mathbf{u}}, \mathbf{v}) = 0$ daje:

\begin{equation}
    L(\mathbf{u}, \mathbf{v}) \approx \sum_j 
    \frac{\partial L(\mathbf{u}, \mathbf{v})}{\partial u_j} 
    (u_j - \tilde{u}_j)
\end{equation}

Doprinos ulaznog neurona $j$ ovoj aproksimaciji definiše se kao njegova relevantnost. Nakon normalizacije ukupnim doprinosom radi ispunjenja uslova konzervacije, relevantnost ulaznog neurona $j$ određuje se udelom u raspodeli relevantnosti izlaznog neurona $i$:

\begin{equation}
    R_j = \sum_i \frac{u_j \frac{\partial L(\mathbf{u}, \mathbf{v})}{\partial u_j}}
    {\sum_k u_k \frac{\partial L(\mathbf{u}, \mathbf{v})}{\partial u_k}} R_i
    \label{eq:deep_taylor}
\end{equation}

pri čemu je uzeto $\tilde{u}_j = 0$ kao prirodan izbor korene tačke, što 
odgovara pretpostavci da neuron bez ulaznog signala ne doprinosi izlazu. 
Detaljano izvodjenje dato je u Dodatku \ref{appendix:deep_taylor}.
Može se pokazati da ova formulacija striktno zadržava svojstvo konzervacije relevantnosti između slojeva, a formalni dokaz ovog održanja naveden je u Dodatku \ref{appendix:conservation}.

Važno je napomenuti da izbor korene tačke $\tilde{\mathbf{u}}$ nije jedinstven 
— različiti izbori vode ka različitim pravilima propagacije i različitim 
vrstama dekompozicije. Formula \eqref{eq:deep_taylor} predstavlja opšti okvir, 
a konkretna pravila ($z^+$, $\alpha\beta$ i druga) mogu se posmatrati kao 
specijalni slučajevi koji nastaju različitim izborima $\tilde{\mathbf{u}}$ 
\cite{montavon2017explaining}.

Interesantno je da se može pokazati da uniformna primena formule 
\eqref{eq:deep_taylor} na sve slojeve mreže rezultuje objašnjenjem ekvivalentnim 
metodi \textit{Input $\times$ Gradient} \cite{montavon2019lrp}. Ovo ukazuje da 
LRP zapravo predstavlja uopštenje gradijentnih metoda, pri čemu pažljiv izbor 
pravila propagacije po slojevima omogućava kvalitativno bolja objašnjenja od 
prostog množenja gradijenata sa ulazom.

\subsubsection{Propagacija kroz linearni sloj}

Za linearni sloj sa ulaznim vektorom $\mathbf{x}$ i matricom težina $\mathbf{W}$, 
aktivacija izlaznog neurona definisana je kao $z_i = \sum_j x_j w_{ji}$. 
Relevantnost ulaznog neurona $j$ računa se kao:

\begin{equation}
    R_j^{(l)} = \sum_i \frac{z_{ji}}{z_i} R_i^{(l+1)}, \quad z_{ji} = x_j w_{ji}
    \label{eq:lrp_basic}
\end{equation}

Neuron $j$ prima udeo relevantnosti od neurona $i$ proporcionalan njegovom 
relativnom doprinosu aktivaciji tog neurona. Međutim, ukoliko je $z_i$ blisko 
nuli usled međusobnog poništavanja pozitivnih i negativnih doprinosa, dolazi do 
numeričke nestabilnosti — što direktno motiviše pravilo propagacije opisano u nastavku.

\paragraph{$z^+$-pravilo.} Originalno je LRP razvijen za konvolucione mreže sa 
ReLU aktivacijama, gde su sve aktivacije nenegativne. U tom slučaju numerička 
nestabilnost se rešava propagiranjem isključivo pozitivnih doprinosa težina 
($w_{ji} > 0$). Međutim, u Transformer arhitekturama aktivacione funkcije 
poput GELU mogu imati negativne vrednosti, pa doprinos $z_{ji} = x_j w_{ji}$ 
može biti pozitivan čak i kada je $w_{ji} < 0$ (ukoliko je $x_j < 0$). Stoga 
se pozitivnost definiše na osnovu predznaka samog proizvoda:

\begin{equation}
    z_{ji}^+ = \max(0, x_j w_{ji})
\end{equation}

a pravilo propagacije postaje:

\begin{equation}
    R_j^{(l)} = \sum_i \frac{z_{ji}^+}{\sum_{j'} z_{j'i}^+} R_i^{(l+1)}
    \label{eq:zplus}
\end{equation}

Posledica ovog pravila je da sve propagirane relevantnosti ostaju nenegativne 
kroz sve slojeve mreže — dokaz je dat u Dodatku \ref{appendix:positive}. 
Postoji uopštenje ovog pravila — $\alpha\beta$-pravilo — koje uzima u obzir 
i negativne doprinose kroz poseban parametar $\beta$, ali ono neće biti 
razmatrano u ovom radu.

\subsubsection{Propagacija kroz neparametarske operacije}

ViT arhitektura sadrži dva tipa neparametarskih operacija koje zahtevaju poseban tretman: \textbf{sabiranje} (rezidualne veze) i \textbf{množenje} (mehanizam pažnje).

\paragraph{Sabiranje.} Za operaciju $\mathbf{z} = \mathbf{u} + \mathbf{v}$, primenom opšte DTD formule dobija se raspodela:

\begin{equation}
    \label{eq:addition}
    R_j^{\mathbf{u}} = \frac{u_j}{u_j + v_j} R_j^{\mathbf{z}}, \quad 
    R_j^{\mathbf{v}} = \frac{v_j}{u_j + v_j} R_j^{\mathbf{z}}
\end{equation}

Iako sabiranje po definiciji očuvava relevantnost (dokaz u Dodatku \ref{appendix:addition}), ono unosi rizik od numeričke nestabilnosti. Ukoliko se komponente ulaza međusobno poništavaju ($u_j + v_j \approx 0$), imenilac teži nuli, što dovodi do deljenje vrednošću bliskom nuli i nekontrolisane eksplozije relevantnosti.

\paragraph{Množenje.} Za matrično množenje $\mathbf{Z} = \mathbf{U}\mathbf{V}$, primenom DTD formule na oba ulazna toka dolazi do dupliranja ukupne relevantnosti na ulazu:

\begin{equation}
    \sum_{i,k} R_{ik}^{\mathbf{U}} + \sum_{k,j} R_{kj}^{\mathbf{V}} = 2 \sum_{i,j} R_{ij}^{\mathbf{Z}}
\end{equation}

što direktno narušava aksiom konzervacije (dokaz u Dodatku \ref{appendix:multiplication}).

\paragraph{Normalizacija relevantnosti.} Kako bi se istovremeno prevazišao problem numeričke nestabilnosti kod sabiranja i narušavanje konzervacije kod množenja, uvodi se postupak stabilizacije i normalizacije relevantnosti \cite{chefer2021transformer}. Za bilo koju binarnu neparametarsku operaciju sa ulaznim tenzorima relevantnosti $\mathbf{R}^{\mathbf{A}}$ i $\mathbf{R}^{\mathbf{B}}$, normalizovana vrednost $\hat{\mathbf{R}}^{\mathbf{A}}$ definiše se kao:

\begin{equation}
    \hat{R}_j^{\mathbf{A}} = R_j^{\mathbf{A}} \cdot \frac{\sum_l |R_l^{\mathbf{A}}|}{\sum_l |R_l^{\mathbf{A}}| + \sum_k |R_k^{\mathbf{B}}|} \cdot \frac{\sum_i R_i^{\text{izlaz}}}{\sum_l R_l^{\mathbf{A}}}
    \label{eq:normalized}
\end{equation}

Ovim izrazom se vrši preraspodela srazmerno doprinosima pojedinačnih ulaza. Primenom apsolutnih vrednosti u imeniocu drugog člana sprečava se deljenje vrednošću bliskom nuli kod sabiranja, čime se vrednost relevantnosti zadržava u zadatom okviru $0 \leq \hat{R}_i^{\mathbf{A}} \leq \sum_j R_j^{\text{izlaz}}$. Istovremeno, skaliranje trećim članom poništava višak relevantnosti nastao množenjem i uspostavlja striktnu konzervaciju (dokaz u Dodatku \ref{appendix:normalization}).

\paragraph{Softmax i LayerNorm.} Za slojeve poput Softmax i Layer Normalization, relevantnost se propagira nepromenjeno. Ovi slojevi se u LRP propagaciji tretiraju kao identičko preslikavanje (\textit{identity mapping}) \cite{chefer2021transformer}.

\subsection{Metode zasnovane na mehanizmu pažnje}

Metode opisane u prethodnoj sekciji oslanjaju se isključivo na gradijente i 
relevantnosti koje se propagiraju kroz težinske matrice. Međutim, 
Transformer arhitekture poseduju jedinstvenu strukturu — mehanizam 
pažnje — koja eksplicitno modeluje interakcije između tokena. Prirodno je 
stoga pokušati iskoristiti ovu strukturu direktno za interpretabilnost.

\subsubsection{Attention Rollout}

Na početku mreže svaki token ima jasan identitet — direktno odgovara određenom 
delu ulazne slike. Međutim, kako signal prolazi kroz slojeve Transformer 
enkodera, tokeni međusobno razmenjuju informacije isključivo putem mehanizma 
višestruke pažnje, što dovodi do postepenog mešanja njihovih reprezentacija. 
Nakon nekoliko slojeva, reprezentacija svakog tokena sadrži informacije iz 
potencijalno svih delova slike, pa sirove mape pažnje jednog sloja ne mogu 
sa sigurnošću reći koji token nosi koju informaciju.

Pored toga, istraživanja su pokazala da dublje slojeve Transformer
mreža prati fenomen kolapsa ranga (\textit{rank collapse}) matrice pažnje — 
u višim slojevima mreže težine pažnje teže ka uniformnoj raspodeli, što znači 
da CLS token posvećuje gotovo jednaku pažnju svim patch tokenima 
\cite{dong2021attention}. Posledica je da sirove mape pažnje iz viših slojeva 
nose sve manje lokalizovanu informaciju o tome koji delovi slike su važni.

Kako bi se pratilo kako se informacija iz pojedinih pečeva propagira 
kroz sve slojeve do CLS tokena, predložen je Attention Rollout 
\cite{abnar2020quantifying}. Ključna opservacija je da je zbir redova matrice 
pažnje jednak jedinici usled softmax normalizacije, što znači da množenje dve 
takve matrice takođe čuva ovo svojstvo. Time se propagacija pažnje kroz više 
slojeva može prirodno modelovati matričnim množenjem.

Zbog rezidualnih veza koje direktno prenose informaciju prethodnog sloja, 
svaki token prima informacije iz dva izvora: kroz mehanizam pažnje i direktno 
od sebe samog. Ovo se modeluje dodavanjem jedinične matrice $\mathbf{I}$, 
pri čemu se vrši normalizacija kako bi se očuvalo svojstvo da zbir redova 
ostane jednak jedinici:

\begin{equation}
    \hat{\mathbf{A}}^{(l)} = \frac{\mathbf{A}^{(l)} + \mathbf{I}}{2}
\end{equation}

Rollout mapa dobija se rekurzivnim množenjem normalizovanih matrica pažnje 
kroz sve slojeve:

\begin{equation}
    \mathbf{C}^{(l)} = \bar{A}^{(1)} \cdot \bar{A}^{(2)} \cdots \bar{A}^{(N)}
\end{equation}

Konačna mapa značajnosti dobija se iz reda koji odgovara CLS tokenu u 
matrici $\mathbf{C}^{(L)}$:

\begin{equation}
    S_{rollout}(\mathbf{x}) = \mathbf{C}^{(L)}_{[\text{CLS}, 1:N]}
\end{equation}

Budući da svaki sloj poseduje $h$ komponenti pažnje, standardni pristup je 
usrednjavanje matrica pažnje po svim komponentama pre množenja:

\begin{equation}
    \mathbf{A}^{(l)} = \frac{1}{h} \sum_{h'=1}^{h} \mathbf{A}^{(l,h')}
\end{equation}

Međutim, kako različite glave pažnje specijalizuju se za različite vrste 
zavisnosti \cite{voita2019analyzing}, jednostavnim usrednjavanjem gube se 
ove specijalizovane informacije i sve glave se tretiraju kao podjednako 
informativne.

Pored toga, Attention Rollout je osetljiv na problem zagušenja 
pažnje (\textit{attention saturation}) — ukoliko jedan token dominira 
matricom pažnje sa izrazito velikom težinom, ova neravnomernost se 
eksponencijalno pojačava množenjem kroz slojeve, što rezultuje mapama 
značajnosti fokusiranim na mali broj tokena. Konačno, metoda ne uzima u 
obzir gradijentni signal niti doprinose FFN slojeva, već sve težine pažnje 
tretira kao podjednako relevantne za konačnu odluku modela. Posledica je 
da je Attention Rollout potpuno nezavistan od ciljne klase, što predstavlja fundamentalno 
ograničenje za primenu u interpretabilnosti i direktno motiviše metode opisane u nastavku.

\subsubsection{Propagacija relevantnosti kroz mehanizam pažnje \textit{(Transformer Attribution)}}

Iako Attention Rollout predstavlja unapređenje u odnosu na sirove mape pažnje, on i dalje koristi same težine pažnje kao aproksimaciju važnosti tokena. Chefer i saradnici \cite{chefer2021transformer} pokazuju da pažnja sama po sebi nije pouzdana mera relevantnosti, jer opisuje koliko jedan token “gleda” drugi, ali ne i koliko taj odnos doprinosi konačnoj odluci modela, zato uvode propagaciju relevantnosti kroz pažnju pomoću LRP-a, uz signal specifičan za ciljnu klasu.

U osnovnoj varijanti metode, umesto matrica pažnje $\mathbf{A}^{(l)}$ koriste 
se njihove relevantnosti $\mathbf{R}_A^{(l)}$ dobijene LRP propagacijom, pri 
čemu se zadržavaju samo pozitivni doprinosi i 
modeluju rezidualne veze dodavanjem jedinične matrice. Time se dobija propagacija 
nalik na Attention Rollout, ali zasnovana na relevantnostima umesto na sirovim težinama pažnje.

\begin{equation}
    \hat{\mathbf{R}}_A^{(l)} = \mathbf{I} + \mathbb{E}_h\left(\mathbf{R}_A^{(l)}\right)^+
\end{equation}

\begin{equation}
    \mathbf{C} = \hat{\mathbf{R}}_A^{(1)} \cdot \hat{\mathbf{R}}_A^{(2)} 
    \cdots \hat{\mathbf{R}}_A^{(B)}
\end{equation}

gde $\mathbb{E}_h$ označava usrednjavanje po komponentama pažnje, $(\cdot)^+$ 
zadržava samo pozitivne vrednosti u skladu sa $z^+$-pravilom, a dodavanje 
jedinične matrice $\mathbf{I}$ modeluje rezidualne veze.

Autori dalje unapređuju ovu osnovu uvođenjem ponderisane relevantnosti pažnje 
(\textit{weighted attention relevance}). Umesto da se koristi samo relevantnost 
$\mathbf{R}_A^{(l)}$, ona se dodatno ponderiše gradijentima matrice pažnje 
$\nabla \mathbf{A}^{(l)}$ u odnosu na ciljnu klasu:

\begin{equation}
    \bar{\mathbf{A}}^{(l)} = \mathbf{I} + 
    \mathbb{E}_h\left(\nabla \mathbf{A}^{(l)} \odot \mathbf{R}_A^{(l)}\right)^+
\end{equation}

Nakon dodavanja jedinične matrice i ponderisanja, svaki red rezultujuće 
matrice $\bar{\mathbf{A}}^{(l)}$ normalizuje se kako bi zbir ostao jednak 
jedinici:

\begin{equation}
    \bar{\mathbf{A}}^{(l)} \leftarrow \frac{\bar{\mathbf{A}}^{(l)}}
    {\sum_j \bar{A}^{(l)}_{ij}}
\end{equation}

Na ovaj način se očuvava konzervacija relevantnosti kroz slojeve, u skladu sa LRP principom. Bez normalizacije, slojevi sa većim gradijentnim doprinosom mogli bi nesrazmerno da dominiraju pri matričnom množenju kroz blokove, što bi narušilo ravnotežu propagacije.


Nakon normalizacije, konačna mapa relevantnosti dobija se propagacijom kroz 
sve blokove matričnim množenjem, kao Attention Rollout, ali 
sada sa ponderisanim relevantnostima umesto sa sirovim težinama pažnje:

\begin{equation}
    \mathbf{C} = \bar{\mathbf{A}}^{(1)} \cdot \bar{\mathbf{A}}^{(2)} 
    \cdots \bar{\mathbf{A}}^{(N)}
\end{equation}
Intuicija iza ovog ponderisanja je 
da gradijent $\nabla \mathbf{A}^{(l)}$ govori koliko je promena u matrici 
pažnje uticala na izlaz za ciljnu klasu, dok relevantnost $\mathbf{R}_A^{(l)}$ 
govori koliko je sama matrica pažnje bila aktivna. Njihov proizvod stoga 
kombinuje oba signala — zadržavaju se samo oni elementi matrice pažnje koji 
su istovremeno i relevantni i osetljivi na promenu ciljne klase.

Zadržavanje samo pozitivnih vrednosti $(\cdot)^+$ je u skladu sa $z^+$-pravilom 
uvedenim u LRP sekciji — propagiraju se isključivo pozitivni doprinosi, što 
garantuje nenegativnost mapa značajnosti. Konačna mapa značajnosti dobija se 
iz reda koji odgovara CLS tokenu:

\begin{equation}
    S_{chefer}(\mathbf{x}) = \mathbf{C}_{[\text{CLS}, 1:N]}
\end{equation}

Za razliku od metode Attention Rollout, ova metoda je zavisna od ciljne 
klase — gradijenti $\nabla \mathbf{A}^{(l)}$ se računaju u odnosu na 
specifičnu klasu, pa različite klase daju različite mape značajnosti za istu 
ulaznu sliku.


\subsubsection{Gradient-Driven Multi-Head Attention Rollout (GMAR)}
Prethodno opisane metode zasnovane na mehanizmu pažnje agregiraju doprinose pojedinih komponenti pažnje jednostavnim usrednjavanjem, tretirajući sve glave kao podjednako
važne. Međutim, kao što je napomenuto u poglavlju o ViT arhitekturi, različite 
komponente specijalizuju se za različite vrste zavisnosti \cite{voita2019analyzing}, 
pa je opravdano pretpostaviti da nemaju isti doprinos konačnoj odluci modela za 
datu klasu. \textit{Gradient-Driven Multi-Head Attention Rollout} (GMAR) 
\cite{jo2025gmar} adresira upravo ovaj problem uvođenjem gradijentno izvedenih težina 
za svaku komponentu pažnje.

Za svaki sloj $l$ i svaku komponentu $h$, važnost komponente određuje se na osnovu 
gradijenata izlaza modela u odnosu na matricu pažnje te komponente:

\begin{equation}
    G_h^{(l)} = \sum_{i,h} \left| \nabla_{\mathbf{A}_{h,ij}^{(l)}} f_c \right|
\end{equation}

Težina svake komponente dobija se normalizacijom ovih gradijenata kako bi zbir 
doprinosa kroz sve komponente bio konzistentan:

\begin{equation}
    w_h^{(l)} = \frac{G_h^{(l)}}{\sum_{h'=1}^{H} G_{h'}^{(l)}}
\end{equation}

Ponderisana matrica pažnje za sloj $l$ tada se definiše kao:

\begin{equation}
    \bar{\mathbf{A}}^{(l)} = \sum_{h=1}^{H} w_h^{(l)} \mathbf{A}_h^{(l)}
\end{equation}

Kako bi se modelovale rezidualne veze, dodaje se jedinična matrica, a zatim se rezultat redno normalizuje:
\begin{equation}
    \bar{A}^{(l)} = \bar{\mathbf{A}}^{(l)} + \mathbf{I}
\end{equation}
\begin{equation}
    \bar{A}^{(l)}_{ij} = \frac{\bar{A}^{(l)}_{ij}}{\sum_j\bar{A}^{(l)}_{ij}} 
\end{equation}

Konačna mapa dobija se rekurzivnim množenjem kroz sve slojeve:

\begin{equation}
    \mathbf{C} = \bar{\mathbf{A}}^{(1)} \cdot \bar{\mathbf{A}}^{(2)} 
    \cdots \bar{\mathbf{A}}^{(B)}
\end{equation}

Na ovaj način GMAR kombinuje prostornu strukturu matrica pažnje sa gradijentnim signalom — glave koje više doprinose 
predikciji ciljne klase dobijaju veću težinu, dok glave sa malim 
gradijentima imaju zanemarljiv uticaj na konačnu mapu značajnosti.

Konačna mapa značajnosti dobija se iz reda koji odgovara CLS tokenu:

\begin{equation}
    S_{GMAR}(\mathbf{x}) = \mathbf{R}^{(L)}_{[\text{CLS}, 1:N]}
\end{equation}

%\paragraph{Prošireni GMAR: GMAR + LRP-relevantnost}

%Kao proširenje, moguće je kombinovati GMAR sa pristupom Chefer-a i saradnika i umesto sirovih matrica pažnje $\mathbf{A}_h^{(l)}$ koristiti relevantnosti matrica pažnje 
%$\mathbf{R}_{A,h}^{(l)}$
% dobijene LRP propagacijom. U tom slučaju ponderisana matrica pažnje za sloj 
%$l$
%l definiše se kao:
%\begin{equation}
%    \bar{\mathbf{A}}^{(l)} = \sum_{h=1}^{H} w_h^{(l)} \mathbf{R}_{A,h}^{(l)}
%\end{equation}


%Ovaj hibridni pristup kombinuje gradijentno ponderisanje komponentata (GMAR) sa relevantnošću iz LRP-a, čime se uvodi dodatni signal relevantnosti uz već prisutan gradijentni signal.